\chapter{Semi- and Self-Supervised Learning}
\label{ch:semisup}
Labels are expensive; raw data are cheap. Placement interviews at product companies increasingly ask
``you have 1000 labelled and a million unlabelled examples --- what do you do?'' The honest answer
starts with the assumptions that make unlabelled data useful, and with the fact that it can hurt.

\section{Assumptions}
\prototype{Two half-moons, six labelled points, six hundred unlabelled ones.}
Unlabelled points tell you about $p(\x)$. They help predict $y$ only if $p(\x)$ carries information about
$p(y\mid\x)$. Three standard ways of saying so:
\begin{itemize}
	\ii \vocab{Smoothness}: points close in a high-density region should share a label.
	\ii \vocab{Cluster} (low-density separation): the decision boundary should pass through low-density
	regions, so each cluster tends to be one class.
	\ii \vocab{Manifold}: data lie near a low-dimensional manifold, and distances \emph{along} it are the
	meaningful ones.
\end{itemize}
When these hold, unlabelled data reshape the boundary; when they fail, they mislead it.

\section{Self-training}
\prototype{Train on the labelled set, label the unlabelled points you are most sure of, retrain, repeat.}
\vocab{Self-training} (pseudo-labelling): fit a model on labelled data; predict the unlabelled data; add
predictions above a confidence threshold as \vocab{pseudo-labels}; refit; repeat. It is simple and
model-agnostic (scikit-learn's \code{SelfTrainingClassifier}). The failure mode is
\vocab{confirmation bias}: early mistakes become training labels and are reinforced, especially when the
initial model is poor or poorly calibrated. Mitigations: high thresholds, class-balanced selection,
re-labelling every round (not freezing pseudo-labels), and the teacher--student variants below.

\section{Co-training and label propagation}
\prototype{A web page described by its text and, independently, by the anchor text of links pointing to it.}

\vocab{Co-training} \cite{cotraining} needs two \emph{views} of each example, each sufficient on its own and
conditionally independent given the label. Two classifiers, one per view, label unlabelled examples for each
other: an example confidently labelled by the text model teaches the link model, which sees it differently.

\vocab{Label propagation} builds a graph with edge weights $W_{ij}=\exp(-\gamma\norm{\x_i-\x_j}^2)$ and lets
labels flow along it. With labelled nodes clamped, the fixed point is the \emph{harmonic} solution: every
unlabelled node's label distribution is the weighted average of its neighbours',
\[ \bm F_u=(\mathbf D_{uu}-\W_{uu})\inv\W_{ul}\Y_l,\qquad \mathbf D=\diag\Bigl(\sum_jW_{ij}\Bigr). \]
\codefile[code/ml/c12\_semisup.py: label propagation in closed form]{gen/ml/snip/c12_semisup-harmonic.py}
On 120 moon points with four labels this matches the iterative \code{LabelPropagation} of
scikit-learn to $10^{-4}$ and labels $\gv{c12_semisup/lp_acc}$ of the points correctly. \vocab{Label spreading} adds a
regularisation $\alpha$ that lets clamped labels be partly overwritten, which tolerates label noise.

\section{Consistency regularisation}
\prototype{An image and its slightly cropped, recoloured copy should get the same prediction.}
Modern semi-supervised deep learning adds an unsupervised loss that penalises a model whose predictions
change under perturbations of an \emph{unlabelled} input --- a direct implementation of the smoothness
assumption. Examples: the $\Pi$-model (two stochastic passes should agree), Mean Teacher (a student matches
an exponential moving average of its own weights), and FixMatch (pseudo-label a weakly augmented view only if
confident, train the strongly augmented view to predict it). The total loss is supervised cross-entropy
plus $\lambda\times$ consistency loss, with $\lambda$ ramped up as the model improves.

\section{When unlabelled data hurts}
\prototype{One cloud of points whose labels flip across a line through its densest region.}

\begin{example}[Helps and hurts, measured]
	Six labelled points (three per class), 594 unlabelled, mean test accuracy over 20 random draws:
	\begin{center}
		\small
		\begin{tabular}{p{0.3\linewidth}cccc}
			\toprule
			data & LR & RBF-SVM & self-training & label spreading \\
			& \multicolumn{2}{c}{(6 labels only)} & \multicolumn{2}{c}{(6 labels + 594 unlabelled)} \\ \midrule
			\gv{c12_semisup/ss_rows}
			\bottomrule
		\end{tabular}
	\end{center}
	On two moons the classes are separated by a low-density gap and label spreading lifts accuracy from
	$\gv{c12_semisup/moon_sup}$ to $\gv{c12_semisup/moon_spread}$. On the second dataset the clusters run
	along $x_1$ but the label depends on the sign of $x_2$: the cluster assumption is false, and label
	spreading ($\gv{c12_semisup/cloud_spread}$) loses to a plain logistic regression trained on the six
	labelled points ($\gv{c12_semisup/cloud_lr}$), whose linear bias happens to be right. Self-training with an
	RBF-SVM made things worse in \emph{both} cases ($\gv{c12_semisup/moon_self}$ on moons): with six labels the
	first model is weak, and its confident mistakes are fed back as truth.
\end{example}

\begin{remark}
	\trap ``More data never hurts.'' More \emph{labelled} data from the same distribution essentially never
	hurts; \emph{unlabelled} data can, when the method's assumption about $p(\x)$ and $p(y\mid\x)$ is wrong,
	or through confirmation bias. Always compare against the supervised baseline on a labelled validation
	set --- which, with few labels, should be cross-validated.
\end{remark}

\section{Contrastive learning in one page}
\prototype{Two augmented views of the same image should be closer to each other than to any other image in
	the batch.}
Self-supervised \vocab{contrastive learning} learns an encoder $f$ without labels: for each example build
two views (augmentations); embeddings of the same example are \emph{positives}, all other pairs in the batch
are \emph{negatives}. The \vocab{InfoNCE} loss for anchor $i$ with positive $i^+$ is
\[ \ell_i=-\log\frac{\exp(\mathrm{sim}(\z_i,\z_{i^+})/\tau)}{\sum_j\exp(\mathrm{sim}(\z_i,\z_j)/\tau)}, \]
a softmax cross-entropy whose ``correct class'' is the positive; $\mathrm{sim}$ is cosine similarity and
$\tau$ a temperature (small $\tau$ sharpens the distribution and focuses on hard negatives).

\begin{example}[InfoNCE on three items]
	Normalised view-1 and view-2 embeddings for three items give, at $\tau=0.5$, the similarity matrix
	$\begin{psmallmatrix}\gv{c12_semisup/nce_S}\end{psmallmatrix}$ (rows: anchors; diagonal: positives).
	Row 1's softmax puts $\gv{c12_semisup/nce_p00}$ on its positive; the mean loss over the three rows is
	$\gv{c12_semisup/nce_loss}$, identical to \code{torch.nn.functional.cross_entropy} applied to the matrix
	with targets $(0,1,2)$.
\end{example}
After pre-training, a linear classifier on the frozen embeddings (``linear probe'') or light fine-tuning needs
far fewer labels. What the augmentations leave unchanged is what the representation learns to ignore, so
augmentation design \emph{is} the inductive bias. The same recipe pairs images with captions in
image--text models; see \otherbook{Deep Learning} for SimCLR- and CLIP-style training.

\begin{moral}
	Unlabelled data help only through an assumption linking $p(\x)$ to $p(y\mid\x)$; check it against a
	supervised baseline, and beware of a model grading its own homework.
\end{moral}

\section\problemhead

\begin{problem}[Pick the method]
	\twochili
	For each scenario name the assumption you rely on and a method: (a) 200 labelled and 50\,000 unlabelled
	product photos; (b) a citation graph with 1\% of papers labelled by field; (c) e-mails described by body
	text and by sender metadata.
	\begin{hint}
		Graphs suggest propagation; two views suggest co-training.
	\end{hint}
	\begin{sol}
		(a) Smoothness and cluster assumptions: consistency regularisation (FixMatch-style), or
		self-supervised pre-training followed by a linear probe. (b) Homophily (linked papers share fields, a graph version of smoothness): label
		propagation or a graph neural network. (c) Two roughly independent views: co-training.
	\end{sol}
\end{problem}

\begin{problem}[Harmonic label]
	\onechili
	An unlabelled node is connected to labelled neighbours with weights $2$ (class A), $1$ (class A) and $3$
	(class B) and to nothing else. After label propagation, what is its probability of class A?
	\begin{hint}
		It takes the weighted average of its neighbours' label distributions.
	\end{hint}
	\begin{sol}
		$(2+1)/(2+1+3)=0.5$: a tie, despite having more A-neighbours, because the single B-edge is strong.
	\end{sol}
\end{problem}

\begin{problem}[InfoNCE temperature]
	\twochili
	An anchor has cosine similarity $0.8$ with its positive and $0.5$ with each of 3 negatives. Compute the
	InfoNCE loss at $\tau=1$ and $\tau=0.1$. Which temperature makes the loss more sensitive to the negatives?
	\begin{hint}
		Loss $=-\log\frac{e^{0.8/\tau}}{e^{0.8/\tau}+3e^{0.5/\tau}}$.
	\end{hint}
	\begin{sol}
		$\tau=1$: $\gv{c12_problems/p12c_1}$. $\tau=0.1$: $\gv{c12_problems/p12c_01}$. A small temperature turns a
		0.3 similarity gap into a large logit gap, so the loss approaches $0$ when the positive wins clearly and
		grows quickly when a negative comes close: low $\tau$ concentrates the gradient on hard negatives.
	\end{sol}
\end{problem}
